\begin{prob}{001}{\hosi 1}{Daisu gakukon テンプレート問題}
$\dfrac{n!}{m^{n-1}}$が整数であるような2以上の整数$m,n$の組をすべて求めよ。
\end{prob}

$\dfrac{n!}{m^{n-1}}$が整数であるとき, $m$がある素数$p$で割り切れるとすれば, $\dfrac{n!}{p^{n-1}}$も整数である。\par
$p\geq 3$であるとする(つまり, $m$がある奇素因数をもつとする)。このとき,
\[\disp\sum_{i=1}^{\infty} \llf \dfrac{n}{p^{i}} \rrf < \disp\sum_{i=1}^{\infty}\dfrac{n}{3^{i}} = \dfrac{n}{2} \leq n-1\]
となる。最後は$n\geq 2$を用いた。最左辺は$n!$が$p$で割れる最大の回数を表しており, それに対して$m^{n}$は$p$で少なくとも$n-1$回割り切れるので, この場合は整数とならない。したがって$m$は$2$の冪であることが必要であるから, $m=2^{s}$ ($s\geq 1$)と書ける。\par
$s\geq 2$であるとする。$m^{n-1}$は$2$で$k(n-1)$回割り切れるが,
\[\disp\sum_{i=1}^{\infty} \llf \dfrac{n}{2^{i}} \rrf < \disp\sum_{i=1}^{\infty} \dfrac{n}{2^{i}} = n \leq 2(n-1)\leq s(n-1)\]
であるから, $n!$より$m^{n-1}$のほうが$2$で割り切れる回数が多く, 不適である。よって$m=2$と決まり, $\dfrac{n!}{2^{n-1}}$ が整数となる$n$を決定すればよい。\par
$n=2^{k} + n_0$ ($0\leq n_0 < 2^{k}$, $k\geq 1$) と書ける。$I=\disp\sum_{i=1}^{\infty} \llf \dfrac{n}{2^{i}}\rrf$とおき, $I=n-1$であるとする。
\beqn
I &=& \disp\sum_{i=1}^{k} \llf \dfrac{n_0}{2^{i}}\rrf= \disp\sum_{i=1}^{k} \parena{2^{k-i} + \llf \dfrac{n_0}{2^{i}} \rrf } = (2^{k} - 1) + \disp\sum_{i=1}^{k} \llf \dfrac{n_0}{2^{i}} \rrf\\
&=& (n-n_0 - 1) + \disp\sum_{i=1}^{k}\llf \dfrac{n_0}{2^{i}} \rrf = (n-n_0 -1) + \disp\sum_{i=1}^{k-1} \llf \dfrac{n_0}{2^{i}} \rrf
\eeqn
により, 整理すると
\[n_0 = \disp\sum_{ i=1 }^{k-1} \llf \dfrac{n_0}{2^{i}} \rrf\]
である。$n_0 = 0$ならばこれは正しい。$n_0 \geq 1$ならば, 先と同様に
\[\disp\sum_{i=1}^{\infty} \llf \dfrac{n_0}{2^{i}}\rrf < \disp\sum_{i=1}^{\infty} \dfrac{n_0}{2^{i}} = n_0\]
であるから不適である。よって$n_0 = 0$に限り, $n=2^{k}$に限る。そして, $n=2^{k}$ $(k\geq 1)$のときを改めて調べると, $I=\disp\sum_{i=1}^{k} 2^{k-i} = 2^{t}-1 = n-1$となるのでよい。よって題意は示された。\qed\\
\\
(\tf{解答2})  $n$が$2$の冪であることの部分のみ示す。
$I=\disp\sum_{i=1}^{M} \llf\dfrac{n}{2^i} \rrf = n-1$となるのが$n$が2冪の場合であることをみたいので, $n=2^{t}a$ ($t\geq 0$, $a$は奇数) とおいたときに$a=1$であることを示せばよい。このとき
\beqn
I &=& \disp\sum_{i=1}^{M} \llf \dfrac{2^{t}a}{2^i}\rrf = 2^{t-1}a+\cdots + 2^{1}a + 2^{0}a + \disp\sum_{i=1}^{M-t}\llf\dfrac{a}{2^{i}}\rrf\\
&=& a(2^{t}-1) + \disp\sum_{i=1}^{M-t}\llf\dfrac{a}{2^{i}}\rrf\\
&=& n-1 = 2^{t}a - 1
\eeqn
は,  $\disp\sum_{i=1}^{M-t}\llf\dfrac{a}{2^{i}}\rrf = a-1$ に同値。（つまり$t=0$の場合を考えることに帰着する）$M-t=N$とおくと, $N$は$2^N\leq a$を満たす最大の整数である。\par
$a>1$であるとしよう。このとき, $N\geq 1$となる。等式$\disp\sum_{i=1}^{N}\llf\dfrac{a}{2^{i}}\rrf=a-1$の左辺は, $a!$を素因数分解したときの$2$の指数である。$a$は奇数であるから, $(a-1)!$も同じ$2$の指数を持つ。したがって, $\disp\sum_{i=1}^{N}\llf\dfrac{a-1}{2^{i}}\rrf=a-1$となり, $\dfrac{(a-1)!}{2^{(a-1)}}$が整数となることを示している。$a-1\geq 2$であるから, これは(1)で示したことに矛盾する。よって$a=1$であり, $n$は$2$のべき乗であることが必要である。

